Tento souhrnný seznam uvádí vybrané zdroje a univerzitní guideliney užitečné pro další studium a adaptaci postupů do praxe na VUT. U zdrojů, které jsou přímo citovány z interního materiálu „AI – 101 tipů“, uvádíme poznámky a doporučené využití v kurzech.

\begin{itemize}
  \item Microsoft 365 / Office 365 a Copilot — přehled nástrojů (Chat, integrované funkce v Word/Excel/PowerPoint) a poznámka, že školní varianty (A1 / A3/A5) obsahují specifické edukační funkce; Copilot Chat je dostupný ve školních verzích a Copilot může rozšířit běžné aplikace o generativní funkce \cite{kb_ai_101_tipu_pdf_04e4a115_page_3_1}.
    \begin{itemize}
      \item Praktické využití: integrace do přípravy materiálů, rychlé sumarizace, generování kvízů a nápovědy pro studenty (využít v pilotu s jasnými pravidly pro akademickou integritu). % (podrobnosti viz repozitář).
    \end{itemize}

  \item Pokrok v matematice (součást Vzdělávacích akcelerátorů Microsoft) — nástroj pro generování a správu procvičovacích úloh, sběr postupů žáků a analýzu silných / slabých stránek ve třídě; v rámci rodiny Microsoft nástroje doplňuje OneNote (převod rukopisu a výpočty) a webová služba math.microsoft.com \cite{kb_ai_101_tipu_pdf_04e4a115_page_8_1}.
    \begin{itemize}
      \item Praktické využití: vhodné pro automatizované cvičení doma i diagnostiku chyb; při nasazení v kurzech Zajistit, aby studenti odevzdávali i postup (procesní hodnocení).
    \end{itemize}

  \item Minecraft Education — edukační světy a „Hour of Code“ aktivity vhodné pro názorné seznámení studentů s principy strojového učení a AI; některé lekce jsou zdarma, širší obsah je dostupný v rámci školních licencí Microsoft 365 A3/A5 \cite{kb_ai_101_tipu_pdf_04e4a115_page_15_2}.
    \begin{itemize}
      \item Praktické využití: hands‑on ukázky pro základní i mezioborové kurzy (viz dostupné světy zaměřené na AI a Hodinu kódu).
    \end{itemize}
\end{itemize}

General background knowledge: Pro inspiraci a implementační patterny je užitečné prostudovat guideliney a design patterny z jiných univerzit (např. UCL, TU Delft, Harvard CS50, Georgia Tech) — tyto zdroje a přenositelné příklady jsou shromážděny a anotovány v online repozitáři příručky (viz ai.vut.cz jako living repository). (general background knowledge)